\section{Related Works}
\label{sec:related}
\subsection{Single-view Adjustment Methods}
Early aesthetic optimization works focus on image cropping~\cite{wei2018good,zeng2019reliable,hong2021composing,wang2023image} to reframe an existing photo to improve visual appeal. They predict cropping windows to maximize the aesthetic score but function only as post-processing once an image is captured. Later works extend aesthetic adjustment to virtual camera motion, reusing cropping datasets for learning camera adjustment that emulate crop movements. Su~\etal~\cite{su2021camera} predict in-plane camera shifts for interactive composition refinement, while Li~\etal~\cite{li2025towards} suggest camera rotations through an iterative feedback loop. Liu~\etal~\cite{liu2023beyond} employ outpainting to suggest cropping windows beyond the image border, but their model assumes a fixed camera with enlarged field of view, lacking true parallax changes. More recent efforts attempt to incorporate broader viewpoint variation. Uchida~\etal~\cite{uchida20253d} combine outpainting with single-image 3D reconstruction to enlarge the search space, and Yao~\etal~\cite{yao2025photography} synthesize poor-to-better perspectives via image-to-video generation. However, they rely on hallucinated content and cannot ensure geometric consistency with real scenes.
Lacking awareness of the underlying 3D structure, single-view adjustment methods are restricted to local composition refinement around the original viewpoint. They struggle to discover more aesthetic angles that involve excluding or introducing scene elements, which requires 3D reasoning beyond observed views. In contrast, our method explicitly models spatially varying scene aesthetics, supporting geometry-aware aesthetic reasoning in 3D.

\subsection{3D-exploration Methods}
3D-exploration methods reason directly in 3D environments. They employ RL or search algorithms to actively explore candidate viewpoints in real or simulated scenes. AutoPhoto~\cite{alzayer2021autophoto} uses RL to navigate through 3D to find aesthetic views, while GAIT~\cite{xie2023gait} trains an RL agent to generate indoor tours in simulation. ViewActive~\cite{wu2024viewactive} uses geometric and semantic cues to optimize viewpoint selection for mesh objects. Skartados~\etal~\cite{skartados2024finding} adopt a genetic algorithm to find aesthetic views in a pretrained NeRF scene. However, they rely on dense captures or prebuilt 3D assets that are costly to obtain, and their iterative RL procedures requires physical adjustments in real world. In contrast, our framework benefits from efficient inference of 3D scene aesthetics from sparse observations, and virtually reasons within the aesthetic field. 
% Moreover, they treat the 3D environment as an external source for scoring 2D views, while we embed aesthetic representations intrinsically in 3D, enabling stable and consistent view suggestions.

\subsection{Feature Distillation in 3D Gaussian Splatting}
Recent advances in 3D Gaussian Splatting~\cite{kerbl20233d,mvsplat,depthsplat} have enabled efficient, differentiable rendering and provided a versatile representation for encoding scene geometry and appearance. Building on this, several studies have explored distilling 2D semantic features~\cite{qin2024langsplat,zhou2024feature,li2025semanticsplat,tian2025uniforward} into 3D Gaussian fields for segmentation tasks. These developments suggest the potential of transferring 2D knowledge into 3D representations. However, unlike semantic features that are largely viewpoint-invariant, aesthetic information is inherently viewpoint-dependent and remains largely unexplored in this context. Our work extends this paradigm by distilling aesthetic perception into 3D Gaussian fields to enable viewpoint-aware aesthetic modeling.
